\documentclass[11pt]{article}
\usepackage[margin=1in]{geometry}
\usepackage{amsmath,amssymb,amsthm}
\usepackage{graphicx}
\usepackage[hidelinks]{hyperref}

\newtheorem{theorem}{Theorem}
\newtheorem{lemma}{Lemma}
\newtheorem{remark}{Remark}

\newcommand{\D}{\mathbb D}
\newcommand{\T}{\mathbb T}
\newcommand{\Aut}{\operatorname{Aut}(\D)}
\newcommand{\HH}{\mathcal H}
\newcommand{\BMOA}{\operatorname{BMOA}}
\newcommand{\MD}{M(\mathcal D_\mu)}
\newcommand{\dA}{\,dA}

\title{A finite-measure subcase of the non-radial Hilbert-matrix norm problem}
\author{arXiv automatic math run \texttt{fa\_banach\_001}}
\date{2026-06-18}

\begin{document}
\maketitle

\section*{Source and question}

The source paper is Bellavita--Stylogiannis, \emph{Hilbert matrix operator
Acting between conformally invariant spaces}, arXiv:2408.01060
\cite{BellavitaStylogiannis}.  Its final section records the following
question.

\begin{center}
\includegraphics[width=0.95\linewidth]{figures/open_problem_crop.png}
\end{center}

\noindent In the notation of the source, for a positive Borel measure \(\mu\)
on \(\D\),
\[
 U_\mu(z)=\int_{\D}\log\left|\frac{1-\overline w z}{z-w}\right|^2\,d\mu(w)
\]
and
\[
 \|g\|_{\MD}
 =
 |g(0)|+
 \sup_{\phi\in\Aut}
 \left(\int_{\D}|g'(z)|^2 U_\mu(\phi(z))\dA(z)\right)^{1/2}.
\]
The source proves the exact formula for radial \(\mu\).  The question asks for
the exact norm when \(\mu\) is non-radial.

\section*{Result}

\begin{theorem}
Let \(\mu\) be a finite nonzero positive Borel measure on \(\D\).  Then
\[
 \|\HH\|_{H^\infty\to \MD}
 =
 1+\frac{\pi}{\sqrt2}\sqrt{\mu(\D)}.
\]
In particular, Question 8.1 has this exact answer for every finite non-radial
measure.
\end{theorem}

\section*{Proof intuition}

For finite \(\mu\), the \(M(\mathcal D_\mu)\)-seminorm can be written as an
average of the usual local \(\BMOA\) Green seminorms, after moving the point of
evaluation by automorphisms.  Degenerating automorphisms toward the boundary
send every fixed point of \(\D\) to the same boundary point.  Thus, for the
logarithm \(L(z)=\log(1/(1-z))\), all local seminorms approach the extremal
\(\BMOA\) value \(\pi/\sqrt2\) simultaneously, and finiteness of \(\mu\) allows
dominated convergence through the integral over \(w\).

The upper bound for \(\HH\) is exactly the same estimate used in the source:
\((1-z)(\HH f)'(z)\) is bounded by \(\|f\|_\infty\).  The lower bound is obtained
by testing \(f\equiv 1\).  The function \(\HH(1)\) has the same boundary local
\(\BMOA\) limit as \(L\), so the upper bound is sharp.

\section*{Proof}

Write
\[
G(z,w)=\log\left|\frac{1-\overline w z}{z-w}\right|^2.
\]
For an analytic function \(g\), set
\[
 \mathcal D_a(g)=\int_{\D}|g'(z)|^2G(z,a)\dA(z),\qquad a\in\D.
\]
The Green kernel is invariant under automorphisms:
\[
 G(\phi(z),w)=G(z,\phi^{-1}(w)),\qquad \phi\in\Aut .
\]
Therefore, by Fubini,
\begin{equation}\label{eq:finite-average}
\int_{\D}|g'(z)|^2U_\mu(\phi(z))\dA(z)
=
\int_{\D}\mathcal D_{\phi^{-1}(w)}(g)\,d\mu(w).
\end{equation}

Let
\[
 L(z)=\log\frac1{1-z}.
\]
The classical computation used in the source paper's proof of the
\(H^\infty\to\BMOA\) norm gives
\[
 \sup_{a\in\D}\mathcal D_a(L)^{1/2}=\frac{\pi}{\sqrt2},
 \qquad
 \mathcal D_{\sigma_r(w)}(L)^{1/2}\longrightarrow \frac{\pi}{\sqrt2}
\]
for every fixed \(w\in\D\), as \(r\uparrow 1\), where
\(\sigma_r(z)=(r-z)/(1-rz)\).  Applying \eqref{eq:finite-average} and dominated
convergence gives
\begin{equation}\label{eq:log-finite-seminorm}
\sup_{\phi\in\Aut}\left(\int_{\D}|L'(z)|^2U_\mu(\phi(z))\dA(z)\right)^{1/2}
=\frac{\pi}{\sqrt2}\sqrt{\mu(\D)}.
\end{equation}

For \(f\in H^\infty\), the source proves the identity
\[
 (1-z)(\HH f)'(z)=b_f(z),\qquad b_f\in H^\infty,\quad
 \|b_f\|_\infty\leq \|f\|_\infty .
\]
Since \(\HH f(0)=\int_0^1 f(t)\,dt\), \eqref{eq:log-finite-seminorm} gives
\[
\|\HH f\|_{\MD}
\leq
\left(1+\frac{\pi}{\sqrt2}\sqrt{\mu(\D)}\right)\|f\|_\infty .
\]
Thus the displayed quantity is an upper bound for the operator norm.

It remains to see that it is sharp.  Put
\[
 K(z)=\HH(1)(z)=\frac1z\log\frac1{1-z},
\]
with the removable value \(K(0)=1\).  The same boundary computation used in
Bellavita--Stylogiannis for the lower estimate gives
\[
 \sup_{a\in\D}\mathcal D_a(K)^{1/2}=\frac{\pi}{\sqrt2},
 \qquad
 \mathcal D_{\sigma_r(w)}(K)^{1/2}\longrightarrow \frac{\pi}{\sqrt2}
\]
for every fixed \(w\in\D\) as \(r\uparrow1\).  One way to see the limit is to
use the source decomposition
\[
 zK'(z)=-\frac1z\log\frac1{1-z}+\frac1{1-z};
\]
the first term is negligible in the degenerating Green seminorm, while the
second has the same boundary limit as \(L'\).
Again \eqref{eq:finite-average} and dominated convergence imply
\[
\sup_{\phi\in\Aut}\left(\int_{\D}|K'(z)|^2U_\mu(\phi(z))\dA(z)\right)^{1/2}
=\frac{\pi}{\sqrt2}\sqrt{\mu(\D)}.
\]
Since \(K(0)=1\) and \(\|1\|_{H^\infty}=1\),
\[
 \|\HH\|_{H^\infty\to\MD}
 \geq
 \|\HH(1)\|_{\MD}
 =
 1+\frac{\pi}{\sqrt2}\sqrt{\mu(\D)}.
\]
Together with the upper bound, this proves the theorem.

\section*{Verification notes and limitations}

The proof uses only the source paper's Hilbert-matrix derivative identity and
its standard \(\BMOA\) boundary computation for \(L\) and \(\HH(1)\), plus the
Green-kernel invariance formula.  The finite-measure hypothesis is essential
for the dominated-convergence step as written.  The result does not settle the
remaining infinite non-radial measures in Question 8.1.

The local run indexes had no exact prior packet for arXiv:2408.01060.  A
bounded web search on 2026-06-18 for the arXiv id, the paper title, and phrases
including ``non-radial'', ``Hilbert matrix'', and \(M(\mathcal D_\mu)\) found
the source paper and nearby Hilbert-matrix papers/surveys, but no later exact
answer to Question 8.1.

\section*{Human-review recommendation}

Review as a likely valid partial answer.  The main point to check is the
asserted boundary-local \(\BMOA\) limit for \(K=\HH(1)\), especially that the
negligible logarithmic term remains dominated uniformly enough to justify the
finite-measure dominated-convergence passage.

\begin{thebibliography}{9}
\bibitem{BellavitaStylogiannis}
C.~Bellavita and G.~Stylogiannis,
\emph{Hilbert matrix operator Acting between conformally invariant spaces},
arXiv:2408.01060, 2024.
\end{thebibliography}

\end{document}
